\clearpage
\section{Analiza domeniului de studiu}

Proiectul DIYBMS (Build Your Own Battery Management System) poate fi analizat într-o gamă largă de domenii, inclusiv electronică, programare, design hardware și energetic. Mai exact, aceasta implică utilizarea microcontrolerelor pentru a citi date în timp real despre starea bateriei și a sistemului utilizând senzorii de care dispunem. Sistemul dat este realizat pe un circuit electronic printat proiectat pentru proiectul dat. Sistemul dat mai dispune și de o interfață web prin intermediul căreia utilizatorul poate monitoriza starea bateriilor și face diferite setări pentru fiecare modul.

\subsection{Descrierea domeniului}
În domeniul electronic, proiectul poate fi analizat prin utilizarea componentelor electronice și a circuitelor printate pentru a dezvolta un sistem de monitorizare și management al bateriilor, care să satisfacă nevoile utilizatorului. Proiectarea unui circuit electronic adecvat este crucială pentru performanța sistemului și siguranța bateriilor.

În ceea ce privește programarea, acest proiect implică dezvoltarea unei aplicații software care să poată procesa datele provenite de la senzori și să ofere utilizatorului informații precise și în timp real despre starea bateriilor. În plus, programarea poate fi utilizată pentru a îmbunătăți sistemul prin integrarea unor funcții de control automat, cum ar fi oprirea sistemului în cazul unor condiții de supratensiune sau supracurent.

Interfața web a proiectului DIYBMS are un rol crucial în monitorizarea și controlul sistemului de gestionare a bateriilor. Utilizatorul poate vizualiza datele de monitorizare a bateriilor în timp real și controla sistemul de încărcare, configurând setările de protecție și încărcare prin interfața web. Aceasta poate fi proiectată pentru a afișa informații sub forma unui grafic sau a unei diagrame și poate oferi notificări în timp real pentru evenimente importante. În general, interfața web facilitează o monitorizare și un control eficient, ajutând utilizatorul să ia decizii mai bune privind gestionarea și încărcarea bateriilor, reducând costurile și timpul necesar pentru menținerea acestora într-o stare optimă. 

În domeniul energiei, proiectul poate fi analizat prin dezvoltarea unui sistem care să permită utilizarea eficientă și sigură a bateriilor, reducând astfel costurile și impactul asupra mediului. În plus, integrarea algoritmilor de filtrare, cum ar fi filtrul Kalman poate îmbunătăți performanța sistemului și stabilitatea acestuia.

\subsection{Programarea microcontrolerelor}

Programarea microcontrolerelor reprezintă procesul de dezvoltare a software-ului care controlează și gestionează funcționalitatea acestor dispozitive electronice integrate. Prin intermediul programelor pentru microcontrolere, se pot implementa diverse funcții și algoritmi pentru controlul și monitorizarea dispozitivelor și a interacțiunii cu mediul înconjurător. Aceste programe pot fi scrise în diferite limbaje de programare, cum ar fi C sau C++, și sunt transferate pe microcontroler pentru a fi executate în mod autonom.

Programarea microcontrolerelor permite configurarea și controlul specific al diferitelor dispozitive electronice. Aceasta include gestionarea interfețelor de intrare/ieșire (GPIO), comunicarea cu alte dispozitive prin protocoale precum I2C sau SPI, precum și implementarea algoritmilor de control și logica necesară funcționării sistemului. Programele pentru microcontrolere sunt esențiale în domenii precum automatizarea industrială, Internet of Things (IoT), roboți și multe alte aplicații în care controlul dispozitivelor electronice este necesar.

\subsection{Programele web}

Programele web reprezintă software-ul dezvoltat pentru a fi executat într-un browser web și pentru a oferi interfață și funcționalitate utilizatorilor prin intermediul paginilor web. Aceste programe pot fi scrise în limbaje de programare precum HTML, CSS și JavaScript și sunt utilizate pentru a crea interfețe interactive și dinamice.

Prin intermediul programelor web, se pot dezvolta aplicații complexe și interactive, cum ar fi platforme de socializare, magazine online, aplicații bancare și multe altele. Programele web permit interacțiunea utilizatorilor cu conținutul și serviciile disponibile prin intermediul browserului web, fără a fi nevoie de instalarea unei aplicații separate pe dispozitivul utilizatorului.

\subsection{Găzduirea paginilor web pe microcontroler}

Un microcontroler poate fi configurat pentru a acționa ca un server web și pentru a găzdui pagini web direct pe dispozitivul respectiv. Acest lucru înseamnă că pagina web însăși este stocată și accesată de pe microcontroler, fără a fi necesară o conexiune externă la un server web separat.

Găzduirea paginilor web pe microcontroler poate fi utilă în aplicații în care este necesară o interfață web simplă și autonomă, cum ar fi sistemele de monitorizare și control localizate într-un spațiu restrâns sau sistemele IoT. Prin intermediul paginilor web hostate pe microcontroler, utilizatorii pot accesa și controla funcționalitatea dispozitivului în mod direct și simplu, fără a fi nevoie de alte componente sau infrastructură suplimentară.

\subsection{Domeniul de aplicare}

Proiectul DIYBMS se concentrează pe dezvoltarea unui sistem de gestionare a bateriilor, având în vedere utilizarea surselor regenerabile de energie, precum panouri solare sau turbine eoliene. Scopul principal al acestui sistem este de a oferi o soluție accesibilă și eficientă pentru încărcarea bateriilor folosind surse de energie regenerabilă.

Un aspect important al proiectului este posibilitatea utilizării bateriilor de tip NiMH. Această caracteristică este deosebit de utilă pentru utilizatorii care nu dispun de baterii Li-ion sau care doresc să utilizeze baterii mai accesibile din punct de vedere al costurilor. Sistemul DIYBMS oferă flexibilitatea de a gestiona și monitoriza bateriile NiMH, oferind astfel o alternativă mai economică și mai accesibilă pentru utilizatori.

Prin intermediul acestui proiect, se urmărește democratizarea accesului la sursele regenerabile de energie și facilitarea încărcării eficiente a bateriilor. DIYBMS oferă o soluție practică și convenabilă pentru utilizatorii care doresc să utilizeze surse regenerabile de energie și să beneficieze de o gestionare eficientă a bateriilor în cadrul propriilor proiecte și aplicații.

\subsection{Argumentarea importanței creării unui astfel de sistem}

Crearea unui sistem de gestionare a bateriilor pentru surse regenerabile de energie are o importanță semnificativă în domeniul electronicelor și al energiei regenerabile. Un aspect crucial în utilizarea bateriilor este încărcarea lor corectă, deoarece o încărcare necorespunzătoare poate duce la riscuri grave, precum incendii sau explozii. Prin implementarea unui sistem de monitorizare și management al bateriilor (BMS), se reduce în mod semnificativ acest risc. Deși majoritatea dispozitivelor electronice și vehiculelor electrice sunt dotate cu BMS-uri, sistemul de încărcare a bateriilor regenerabile în gospodăriile individuale nu beneficiază întotdeauna de această tehnologie.

Prin crearea unui sistem DIY de gestionare a bateriilor, așa cum este prezentat în acest proiect, se oferă o alternativă accesibilă pentru încărcarea bateriilor din surse regenerabile. Această soluție poate fi aplicată în diverse domenii, cum ar fi panourile solare, turbinele eoliene sau generatoarele hidraulice, pentru a asigura o încărcare eficientă a bateriilor și pentru a le proteja de deteriorările cauzate de încărcarea inadecvată.

De asemenea, un astfel de proiect poate fi deosebit de util în zonele rurale sau izolate, unde accesul la rețele electrice este limitat. Prin intermediul acestui sistem de gestionare a bateriilor, oamenii pot utiliza bateriile pe care le dețin deja și le pot încărca în mod eficient și sigur, fără a fi nevoie să achiziționeze baterii noi sau să suporte costuri suplimentare pentru dispozitivele BMS pre-fabricate. Astfel, crearea unui astfel de proiect poate aduce beneficii semnificative atât pentru mediul înconjurător, cât și pentru comunitățile umane.

Prin urmare, dezvoltarea și implementarea unui sistem DIY de gestionare a bateriilor pentru surse regenerabile de energie reprezintă o inițiativă importantă, având potențialul de a sprijini utilizarea durabilă a energiei și de a oferi o soluție practică și accesibilă în gestionarea și încărcarea bateriilor.

\subsection{Analiza sistemelor existente}

În cadrul analizei sistemelor existente pentru proiectul de practică tehnologică, s-au identificat trei sisteme de management al bateriilor (BMS - Battery Management System): ANT BMS, Orion BMS și LTT BMS. Toate aceste sisteme își propun să controleze și să monitorizeze performanța bateriilor, având funcționalități similare cu cele prevăzute în proiectul de practică tehnologică.

ANT BMS este un sistem de management al bateriilor dezvoltat de Advanced New Technologies (ANT) și este utilizat în principal în aplicații marine, staționare și pentru vehicule electrice. Acesta monitorizează parametrii cheie ai bateriei, precum tensiunea, curentul și temperatura, pentru a evalua starea de sănătate a acesteia și a preveni supraîncărcarea sau descărcarea excesivă. ANT BMS se remarcă prin flexibilitatea sa, putând fi configurat pentru o gamă variată de aplicații și tipuri de baterii. De asemenea, dispune de funcții de securitate, inclusiv protecție împotriva suprasarcinii și scurtcircuitului.

Orion BMS este un sistem de management al bateriilor dezvoltat de Ewert Energy Systems și este utilizat în special în aplicații de vehicule electrice, inclusiv vehicule cu două roți, mașini, autobuze și vehicule industriale. Acesta monitorizează performanța bateriei și furnizează informații despre tensiune, temperatură și starea bateriei. Orion BMS se remarcă prin performanța sa superioară în ceea ce privește precizia și stabilitatea. De asemenea, oferă o gamă extinsă de funcționalități, inclusiv monitorizarea incendiilor, detectarea scurgerilor de curent și funcții de gestionare a energiei.

LTT BMS este un sistem de management al bateriilor dezvoltat de LTT Global Battery Solutions și este utilizat în principal în aplicații industriale, cum ar fi energia solară și stațiile de încărcare pentru vehicule electrice. Acesta monitorizează performanța bateriei și controlează procesul de încărcare și descărcare pentru a menține bateria în condiții optime de funcționare. LTT BMS se remarcă prin simplitatea și ușurința sa de utilizare. De asemenea, oferă o gamă largă de caracteristici, inclusiv protecție împotriva suprasarcinii și supratensiunii, precum și diverse module de comunicare.

Sistemele analizate (ANT BMS, Orion BMS și LTT BMS) oferă soluții eficiente pentru monitorizarea și controlul performanței bateriilor în diferite domenii de aplicare. Fiecare sistem prezintă avantaje specifice, cum ar fi flexibilitatea, performanța sau ușurința de utilizare, iar acestea pot fi luate în considerare în dezvoltarea proiectului de practică tehnologică.

\subsection{Scopul și obiectivele}

Scopul proiectului este să adaptăm sistemul existent pentru a putea funcționa cu baterii de tip NiMH și să îmbunătățim performanțele sistemului de gestionare a bateriei. Obiectivele specifice ale proiectului sunt următoarele:

\begin{itemize}
\item Proiectarea și dezvoltarea unei soluții pentru compatibilitatea cu bateriile NiMH;
\item Implementarea și testarea soluției proiectate;
\item Optimizarea performanțelor sistemului de gestionare a bateriei;
\item Verificarea și validarea noii soluții de gestionare a bateriei;
\item Documentarea completă a procesului și soluției noi.
\end{itemize}

În general, proiectul are ca scop îmbunătățirea compatibilității și eficienței sistemului de gestionare a bateriei prin adaptarea sa la bateriile de tip NiMH. Scopul final constă în creșterea duratei de viață a bateriilor și reducerea costurilor de întreținere ale sistemului.
